\section{蒙古征服：从联盟网络到区域接收}

一二零六年铁木真获得成吉思汗称号，是蒙古诸部联盟重组的重要节点，却不是一部预定“世界征服史”的起跑线。草原上的亲属、部众、俘虏、贸易与军事关系需要不断维系；对西夏、金和南宋的战争又分别在沙漠边缘、华北城市和长江流域展开。蒙古军的机动、组织与吸纳能力值得解释，但不能把被征服地区的分裂、投附、抵抗、技术人员和地方供给从征服故事中删去。

\subsection{征服改变的是联系，不只是王朝名号}

西夏在一二二七年的覆亡、金在一二三四年的终结，都伴随多年围城、破坏、迁徙和政治重组。旧国家的道路、税源、城镇、工匠和官吏并没有随帝号一起消失；它们成为新的统治者能够利用、又必须重新分配的资源。对华北而言，战争使人口与耕地记录极不稳定，任何前后户数的巨大差异都需要问：是死亡、逃亡、漏报、合并户籍，还是档案口径改变？对河西与西北而言，绿洲、寺院、站道和多语文书同样决定统治能走多远。

南宋的征服尤其不能被写成北方战争的简单延长。四川、两淮、荆襄、江南和海上有不同的地理与供给条件。襄阳、樊城的持久战把江汉水路、造船、粮运与攻城能力推到前台；其后的临安接收和崖山终局，则显示政权崩解发生在不同地方、不同社会群体的不同时间。胜者的军令与遗民的哀歌都保存了真实经验，却各自选择了可以叙述的人和代价。

\subsection{交读窗口：零丁洋上的遗民声音}

一二七九年，文天祥在被元军拘押、经零丁洋北送途中作《过零丁洋》。这是一首以身世、国难和自我誓约相连的七言律诗，末联“人生自古谁无死，留取丹心照汗青”尤其塑造了遗民的道德位置。把它放在襄樊失守、临安接收和崖山之后，并不是要以一位大臣替南方社会发言；它真正照亮的是政权终局如何被一位战败者以诗歌转化为可传述、可争论的忠义记忆。

崖山的日期与战事脉络仍要以《宋史》《元史》并读，并以宋末文集、降附文书、墓志、港口与海战相关遗存校正；《元朝秘史》、蒙古文残存材料和波斯文史籍可以展开更宽的征服网络，却各有译本、成书层次、地名对应与宫廷立场的限制。钱穆式的提问会追究南宋何以不能再维系政治与财政的结合；吕思勉式的提问则要求看见逃亡、降附、海商和军户的不同生计。文天祥的诗能说明一位士大夫选择以何种语言拒绝终局，不能给出投降者人数、普通居民的意愿，或元军在每一处港口的实际行政。

\subsection{征服国家的制度试验}

蒙古统治者并不是以一套不变的草原制度覆盖所有地方。汉地的行中书省和路府、北方的军政组织、站赤系统、地方税源以及各类宗教与商贸中介，构成可变的接收办法。驿站可以加快军令与人员流动，也会把马匹、劳力和供给负担带到沿线家庭；优待商人和跨区域流通可以扩大财政与信息，亦会产生特权、折纳和地方冲突。制度能力由此增长，却从来不是无成本的。

这也是多政权时代留给征服者的遗产。辽的多中心治理、金对华北城市与军户的接收、南宋对江海财赋的经营，并没有被整齐继承，却提供了道路、人口、技术和行政经验。元朝后来能把大都、江南、东北、西北与海港纳入同一帝国尺度，部分正因为它在征服中不断借用和改造这些既有联系；部分限制也来自同样的区域差异，不能只靠一套中枢诏令克服。

\subsection{谁在叙述征服}

《元朝秘史》、汉文《元史》、波斯文与其他区域材料的时间、人物称谓和因果重点并不总是一致。它们不能被熔成一种“蒙古人的声音”，也不该因成书较晚而一概废弃。与城址、墓葬、碑刻、文书和地方记忆相互校验，可以确认若干军事与政治节点，也会使战损、人数、投降动机和具体命令的细节保持不确定。本章保留这种不确定，不是削弱征服的解释，而是拒绝让后来的帝国叙事替所有被卷入战争的人作证。
